\section{Expanded proofs of the generic theorems}
\label{app:proofs}

This appendix writes the central arguments at the level needed to translate
them to the repository's de Bruijn-scoped natural deduction.  We suppress only
routine reassociation of list contexts.

\subsection{One-certificate sealing}

Let \(V\in\Fm\,\PTm\,(n+1)\), \(G\in\Fm\,\PTm\,n\), and
\(\Gamma\subseteq\Fm\,\PTm\,n\).  The constructor
\[
\Prf.\mathsf{exE}:
 \frac{\Gamma\turnstile\exists V
       \qquad V::\Gamma^{\up}\turnstile G^{\up}}
      {\Gamma\turnstile G}
\]
is exactly \cref{thm:seal}.  The side condition traditionally expressed as
``the eigenvariable is not free in \(G,\Gamma\)'' is enforced by the types:
\(G,\Gamma\) live at arity \(n\), while \(V\) lives at arity \(n+1\).
Consequently no separate freshness proof is needed.

For conditional closure, start under context
\((\exists V)::\Gamma\).  The first premise of existential elimination is
\code{Prf.hyp0}.  Its second premise expects
\[
 V::(\exists V)^{\up}::\Gamma^{\up}\turnstile G^{\up}.
\]
Weakening the original extractor inserts \((\exists V)^{\up}\).  Close
existential elimination, then implication introduction.  This is why the
source candidate's proof is only three lines.

\subsection{Dependent sealing with exact shifts}

Let
\[
 A\in\Fm\,\PTm\,(n+1),
 \qquad B\in\Fm\,\PTm\,(n+2),
 \qquad G\in\Fm\,\PTm\,n.
\]
Assume
\[
\begin{array}{ll}
 p:&\Gamma\turnstile\exists A,\\
 q:&A::\Gamma^{\up}\turnstile\exists B,\\
 r:&B::A^{\up}::\Gamma^{\up\up}\turnstile G^{\up\up}.
\end{array}
\]
The second premise for eliminating \(\exists B\) under the opened \(A\) is
exactly \(r\), because shifting the local context
\(A::\Gamma^{\up}\) gives \(A^{\up}::\Gamma^{\up\up}\).  Thus
\[
  q, r\quad\leadsto\quad
  A::\Gamma^{\up}\turnstile G^{\up}.
\]
Eliminate \(\exists A\) with \(p\) to obtain \(G\).  For a chain of length
\(k\), repeat.  For a DAG, choose a nesting order compatible with dependencies;
all consumers of a certificate are proved before its scope closes.

\subsection{Boundary collection}

Suppose a compiled continuation has context
\(B_r::\cdots::B_1::\Gamma\).  Repeated implication introduction gives
\[
 \Gamma\turnstile B_1\to\cdots\to B_r\to G,
\]
with the order reversed relative to the context stack.  If a user-facing order
is required, use weakening and hypotheses to prove the desired permutation,
or compile a conjunction
\(B_1\wedge\cdots\wedge B_r\) and eliminate its fields.  The context form is
canonical for Lean because \code{Prf.trans} takes a family of derivations for
all members.

If a producer \(p_i:\Gamma\turnstile B_i\) later becomes available, apply
\code{Prf.trans} to replace that occurrence while leaving the other boundaries
as hypotheses.  Thus packet completion is monotone: adding producers does not
change downstream extractors.

\subsection{Least accepted certificate}

Let \(z(c,q)=0\) reflect \(V(q,c)\), and abbreviate
\(m=\mu_z(B,q)\).  The checked theorem \code{search\_found} has, schematically,
\[
  z(c,q)=0,	hinspace c<B
   \quad\turnstile\quad
  m<B\wedge z(m,q)=0\wedge m\le c.
\]
Under an existential premise
\(\exists c(c<B\wedge V(q,c))\), open \(c\), convert \(V\) to \(z=0\), and
apply the theorem.  Its conclusion contains no \(c\) except the optional
\(m\le c\) field.  Discard that field before closing the existential, obtaining
\(m<B\wedge z(m,q)=0\).  Reflect back to \(V(q,m)\).

For minimality, \code{search\_minimal} gives
\[
 d<m\quad\turnstile\quad z(d,q)\ne 0.
\]
Reflection yields \(\neg V(q,d)\).  If a first-order bounded universal is
needed, introduce arbitrary \(d\), assume \(d<m\), and derive the implication;
if a qf consumer is enough, retain the reflected no-counterexample term.

\subsection{Trace reflection in arithmetic detail}

Take Boolean failure terms, so every failure is \(0\) or \(1\).  Define
\[
  b_0=0,\qquad b_{i+1}=b_i+e_i.
\]
The checked natural-number arithmetic proves
\[
 b_i+e_i=0\to b_i=0\wedge e_i=0.
\]
One proof uses \(b_i\le b_i+e_i\) and \(e_i\le b_i+e_i\), then equality to
zero and antisymmetry.  Another uses the zero-or-success decomposition: if
either summand were a successor, the sum would be a successor and hence
nonzero.  The first style reuses order lemmas; the second often gives a smaller
proof term.

The induction formula
\[
 \Phi(i)=\bigl(b_i=0\to I(s_i)=1\bigr)
\]
is qf because equality and implication preserve \(\Fm.\qf\).  At zero,
\(s_0=\mathsf{init}\) and \(b_0=0\); congruence transports the initialization
proof.  At a successor, assume \(b_{i+1}=0\), split it, apply the induction
hypothesis to \(b_i=0\), and instantiate local preservation at the current
state and decoded action.  Defining equations for \(s_{i+1}\) complete the
congruence step.

If \(\Trace(N,\tau)\) also contains a terminal check, put that check in
\(e_N\) and run the recursion for \(N+1\), or conjoin it after the prefix
invariant.  The latter keeps the state horizon conventional.

\subsection{Exact observers}

Assume
\[
 V(c),V(d)\turnstile O(c)=O(d).
\]
Define the output formula with an existential certificate.  To prove
single-valuedness, start from two output hypotheses.  Eliminate the first
existential, then the second.  The local context contains
\[
 V(c), O(c)=o, V(d), O(d)=o'.
\]
By independence, \(O(c)=O(d)\).  Equality symmetry and transitivity give
\(o=o'\).  Since the result is in the outer scope, close both existential
eliminations.

The same proof works for an observed relation \(R(c,o)\) if pairwise valid
certificates and their output witnesses imply output equality.  Thus the
observer need not be a term; it may itself be a small relational certificate.

\subsection{Diameter and gap arithmetic}

Suppose a rational metric is represented by a nonnegative numerator after a
common positive denominator is cleared.  The triangle inequality has the
form
\[
 D(x,z)\le D(x,y)+D(y,z).
\]
Under the four assumptions of the gap lemma,
\[
\begin{aligned}
 g&\le D(\ell_i,\ell_j)\\
  &\le D(\ell_i,o)+D(o,o')+D(o',\ell_j)\\
  &\le e+\delta+e<g,
\end{aligned}
\]
a contradiction by transitivity and irreflexivity of strict order.  In object
arithmetic, each rational is a numerator/denominator record and the final
strict inequality is cross-multiplied by the product of positive denominators.
All positivity assumptions belong to the validator.

\subsection{Finite-strategy soundness}

For a \(\forall x<B\exists y<C(x)\) formula, a strategy table is simply a
bounded array \(y_x\).  Its verifier says every digit is below \(C(x)\) and
the matrix holds at \((x,y_x)\).  To prove the formula, introduce arbitrary
\(x\), assume \(x<B\), instantiate the verifier's no-failure theorem at \(x\),
and introduce \(y_x\).  For one more universal level, table entries are
subtables indexed by the preceding history.  Fixed-shape soundness is therefore
a finite nesting of the same argument.

Completeness uses bounded search for a winning response at each existential
history.  A universal node packages all finitely many recursively generated
subtables.  The proof that the package verifies is a qf induction over the
history enumeration.  The table code bound is the product of explicit move
radices.  No infinite strategy or function space occurs.
